\documentclass[a4paper,12pt]{article}

\usepackage[utf8]{inputenc}
\usepackage[T1]{fontenc}
\usepackage[francais]{babel}
\usepackage{mathptmx}
\usepackage[scaled]{helvet}
\usepackage{courier}

\usepackage[margin=2cm,nohead,foot=1.4cm]{geometry}
\usepackage{parskip}
\usepackage{comment}
\usepackage{graphicx}
\usepackage{fancyvrb}

\begin{document}

\begin{flushleft}
\begin{tabular}[t]{l}
\Large\textsf{\textbf{Rapport sur le manuscrit de Thèse présenté par}}\\[4mm]
\Large\textsf{\textbf{Fabrice Dupros}}\\[4mm]
\large\textsf{\textbf{pour l'obtention du Doctorat de l'Université de Bordeaux I}}\\
\end{tabular}%

\vspace{4mm}

\noindent%
\begin{tabular}[t]{l}
Rapporteur: Stéphane \textsc{Genaud},\\
Maître de conférences habilité à diriger les recherches,\\
Université de Strasbourg.\\[5mm]
Titre du document évalué :\\\\[3mm]
\end{tabular}%

\begin{quote}
\Large \textit{\textsf{\textbf{
Contribution à la modélisation numérique de la propagation des ondes sismiques sur architectures multic{\oe}urs et hiérarchiques}}}
\end{quote}

\end{flushleft}

\vspace{3mm}

Le document présenté par Fabrice Dupros décrit les recherches qu'il a menées dans
le cadre de sa thèse sous la direction conjointe de Dimitri Komatitsch, Professeur à l'Université 
de Pau, membre de l'équipe projet INRIA Magique 3D, et Jean Roman, professeur à l'Université de
Bordeaux I, responsable de l'équipe projet INRIA HIEPACS.



\subsection*{Contexte de l'étude}
Le cadre général du travail est celui de la mise en oeuvre efficace 
d'applications de modélisation sismique sur des architectures de 
calcul modernes. Ce travail étudie les moyens d'améliorer l'efficacité
de méthodes numériques classiques, la méthode des éléments finis, et celle
des différences finies fréquement utilisées dans ce type d'application,
pour des processeurs multi-c{\oe}urs disposant d'une hériarchie de mémoires.

Le travail touche un sujet crucial et d'actualité qui concerne toute l'industrie
du logiciel et du matériel : comment programmer efficacement des
ordinateurs dont l'architecture est composée de différents niveaux
de parallélisme (par exemple plusieurs processeurs contenant chacun
de plus en plus de c{\oe}urs), et de mémoires organisées hiérarchiquement.
Si la multiplication des unités de calcul est la seule voie envisageable
aujourd'hui pour croître en puissance de calcul, les exploiter
au mieux est particulièrement difficile. C'est pourtant un défi
que la simulation numérique, l'un des domaines qui peut bénéficier
le plus de puissance de calcul accrue, doit relever.


\subsection*{Analyse du document présenté}

Le document est organisé en deux parties comptant 6 chapitres qui
décrivent le c{\oe}ur du travail.

Un chapitre d'introduction générale, très concis, présente d'abord les enjeux
de la simulation numérique en géophysique et les caractéristiques des  
méthodes numériques proposées dans la littérature. 
Il esquisse ensuite l'évolution des architectures de calcul, multi-c{\oe}urs,
avec des temps d'accès à la mémoire non-uniformes (NUMA).
Suit un premier chapitre (\textit{Cadre de l'étude}) qui reprend en
détails ces éléments. Il décrit plus précisément les méthodes
des différences finies et des éléments finis et le traitement particulier
des bords du domaine (conditions aux limites), puis expose les
évolutions marquantes des architectures de calcul, notamment 
celles concernant l'organisation de la mémoire par rapport aux
unités de calcul. Logiquement, les évolutions des modèles de 
programmation pour exprimer le parallélisme sont introduites 
à la suite. On peut regretter que la discussion de cet état de
l'art ne soit pas plus argumentée concernant les alternatives
non choisies pour le reste de l'étude, comme les modèles de 
programmation de type PGAS. Davantage de références internationales
auraient donné plus de poids aux choix des solutions étudiées
en profondeur par la suite. Ces solutions sont néanmoins parfaitement 
judicieuses car reconnues au meilleur niveau international et
parfois même présentes au sein de logiciels majeurs, comme hwloc,
sous-projet d'OpenMPI.
Ce chapitre pose au final une hypothèse forte du travail, qui
est que le multi-threading est le meilleur modèle de programmation
candidat. Il est invoqué, à juste titre, l'argument que le modèle à
passage de messages (MPI) prépondérant jusqu'à récemment, induit
des surcoûts prohibitifs en mémoire (zones d'échanges). Cependant,
l'idée d'utiliser MPI dans un mode hybride n'est pas écartée,
annonçant la solution retenue au chapitre 5.\\


La première partie traite ensuite de la modélisation élastodynamique.
Elle est constituée de quatre chapitres.  

Le chapitre 2 traite de l'équilibrage de charge en prenant pour motivation
les calculs plus longs sur les bords des domaines (limites absorbantes)
qu'ailleurs. M. Dupros justifie le choix d'un 
découpage quasi-statique permettant de ne pas perdre les avantages liés
à la régularité de la grille, puis il propose un code simplifié
(un Jacobi) présentant une structure de boucles imbriquée similaire au code visé, 
afin d'évaluer diverses implémentations multi-threadées.
Les tests menés comparent le comportement des threads POSIX avec 
GNU OpenMP et des environnements de recherche de l'Université 
de Bordeaux (GOMP, ForestGOMP utilisant les threads \textsc{Marcel}, puis MPC),
sous deux angles : l'opportunité de faire du multi-threading imbriqué
et la capacité à équilibrer dynamiquement la charge (automatiquement
via les directives).
Dans ce cas test faisant abstraction des accès mémoire, l'évaluation 
montre qu'une efficacité quasi idéale peut être atteinte dès qu'on
a au moins un thread par c{\oe}ur, mais que l'imbrication des threads
est mal supportée. Dans ces expériences, la démarche de réduire au minimum
le noyau de calcul pour isoler l'effet observé lors de l'évaluation est 
très pertinente.


%----
Le chapitre 3 (\textit{Prise en compte de l'affinité mémoire}) a pour
objectif de clarifier les effets de la localité des données par 
rapport aux calculs. L'évaluation couvre un nombre conséquent (6) 
d'architectures NUMA différentes. L'expérimentation procède à nouveau 
par tests unitaires pour mettre en évidence les effets NUMA.
L'étude de l'affinité compare l'environnement BubbleSched, permettant
à l'utilisateur de délimiter des ensembles de threads, et à travers
cela, la localité des données à certains ensembles, 
et la politique de placement des données par défaut de Linux. 
Les résultats montrent que des gains importants, de l'ordre
de 10\% (16 c{\oe}urs) à 40\% (512 c{\oe}urs) sont envisageables.


%---
Le chapitre 4 (\textit{Décomposition espace-temps et réduction 
du trafic mémoire}) s'attache au problème des accès mémoires.
L'exposé montre clairement l'aspect irrégulier des accès mémoire
dans la méthode des différences finies. Il est ensuite proposé
de mettre en {\oe}uvre la technique bien connue du \emph{time skewing}
(ré-organisation des accès) sur ce problème. 
L'exposé débouche sur l'énoncé d'un algorithme
général qui découle des contraintes du cas d'espèce.
Bien qu'introduit pédagogiquement, on aurait aimé que celui-ci
soit formalisé.
L'évaluation est approfondie, montrant des résultats convergents, 
sur le noyau Jacobi ainsi que sur le noyau du code sismique,
d'un gain d'au moins 25\% en cycles et défauts de cache. 
Ces résultats sont comparés à la littérature et sont convaincants
quant au fait que le time skewing permet de réduire significativement
l'impact de la hiérarchie mémoire. 
Le deuxième volet de l'expérience ajoute à la stratégie du \emph{time skewing}
la politique de placement mémoire. Ce choix permet de vérifier que 
les deux stratégies n'ont pas d'effets antagonistes. 
Les résultats finaux montrent une amélioration convergente du 
comportement sur la plupart des architectures. L'expérience,
répétée sur différentes architectures, met en évidence 
les différences très importantes de comportement, avec
un gain quasi-nul quand le bus mémoire représente un 
goulot d'étranglement.  
L'expérience met en évidence l'impact du trafic mémoire en fonction
des architectures. 

%----
Le chapitre 5 propose de combiner les modèles de programmation
multi-threadés et ceux à passage de messages dans l'optique d'utiliser
des clusters de n{\oe}uds multic{\oe}urs. Si cette approche dite
\emph{hybride} est bien connue, il n'est pas donné d'indication
sur l'avantage qu'elle aurait en comparaison d'une version 
uniquement multi-threadée.
La comparaison est faite avec une version purement MPI.
Poursuivant le cheminement, une expérience est réalisée avec 
le code de sismique sur un
exemple synthétique en utilisant un cluster complet. 
Les résultats montrent un déséquilibre de charge inférieur pour
la version hybride, permettant un gain en performance de l'ordre
de 25\%.\\


La deuxième partie, beaucoup plus courte, traite de la méthode
des éléments finis. On y trouve deux chapitres.
Le chapitre 6 discute d'abord des problématiques liées aux noyaux 
de calculs matriciels, et en particulier le choix d'un solveur
répondant aux exigences du cas applicatif. Pour cette partie
applicative, il est souligné l'importance du partitionnement
des éléments du maillage. La contribution de ce chapitre réside 
dans un algorithme de coloration de graphe original. 
L'algorithme permet de déterminer l'affectation des éléments de 
maillage aux c{\oe}urs avec l'objectif de minimiser le déséquilibre
en calculs.
Le chapitre 7 présente finalement une simulation numérique en
vraie grandeur en utilisant 1024 c{\oe}urs avec des résultats
convaincants.
 

\subsection*{Conclusion}

Le document est très bien rédigé et agréable à lire. 
Fabrice Dupros y démontre des connaissances poussées à la fois dans le 
domaine des méthodes numériques utilisées pour la géophysique, et
dans celui de la programmation parallèle des architectures 
de calcul moderne. 
Il faut souligner ici la grande variété des connaissances qu'il a
fallu acquérir dans le domaine du parallélisme pour mettre en {\oe}uvre
cette étude. 

La contribution présentée est précieuse et assez rare : il est 
généralement difficile de comparer de manière rigoureuse, dans
un vrai environnement applicatif, plusieurs solutions concurrentes.
Ce travail est particulièrement intéressant aussi car la démarche
adoptée pour le conduire est rigoureuse. Les différents aspects
qui ont un impact sur la performance ont été isolés, étudiés séparément
puis assemblés pour en comprendre les interactions. C'est donc 
une étude expérimentale scientifiquement bien fondée. On peut regretter
cependant qu'une annexe ne liste pas les détails (numéros de versions, 
implémentations MPI) des logiciels utilisés, ceci à des fins de 
reproductibilité des expériences.

Les travaux ont fait l'objet de plusieurs publications (en particulier,
une revue internationale du meilleur niveau, et deux conférences internationales).
Considérant tous ces éléments, j'émets un avis favorable à la soutenance de 
la thèse de Fabrice Dupros en vue de l'obtention du titre de Docteur 
en informatique de l'université de Bordeaux I. La thèse peut être
soutenue en l'état.

\vspace{.8cm}
\begin{flushright}
Stéphane \textsc{Genaud},\\
{\small fait à Strasbourg, le 22 novembre 2010}.\\
%\includegraphics[width=.20\textwidth]{signgenaud.jpg}
\end{flushright}
\end{document}
